% ================= 第 3 章 (前半): 局部-整体相容性猜想 =================

\section{局部-整体相容性猜想}

对权 $k \ge 2$ 的任何本征尖点模形式 $f$, 我们在模曲线的上同调中构造一个可容许的酉 $\GL_2(\Qp)$-Banach 空间 $\Pi_p(f)$.
当与 $f$ 相关联的 $p$ 进 Galois 表示在 $p$ 处潜在晶体且可约时, 我们猜想 $\Pi_p(f)$ 恰好是 \ref{sec:2.3} 的 Banach 空间.

\subsection{\mbox{}}\label{sec:3.1}

我们首先给出关于 \cite{ref:24} 和 \cite{ref:9} 的 $\GL_2(\Qp)$-Banach 空间 $\widehat{H}{}_{c}^{1}(K_1^{p}(M)) \otimes_{\Zp} \Qp$ 和 $\widehat{H}{}^{1}(K_1^{p}(M)) \otimes_{\Zp} \Qp$ 的一些回顾与预备知识.
对 $\GL_2(\mathbb{A}_f)$ 的任何紧开子群 $K_f$, 记:
\[
Y(K_f) := \GL_2(\Q)\backslash \GL_2(\mathbb{A})/\mathbb{C}^{\times}K_f,
\]
其中通过 $a + ib \mapsto \begin{pmatrix} a & -b \\ b & a \end{pmatrix}$ 把 $\mathbb{C}$ 嵌入 $\GL_2(\mathbb{R})$.
对每个与 $p$ 互素的正整数 $M$ 和每个正整数或 $0$ 的整数 $r$, 设:
\[
K_1^{p}(M) := \left\{ \begin{pmatrix} a & b \\ c & d \end{pmatrix} \in \GL_2(\widehat{\mathbb{Z}}{}^{p}),\ (c,d) \equiv (0,1) \bmod M \right\},
\]
\[
K_p(r) := \left\{ \begin{pmatrix} a & b \\ c & d \end{pmatrix} \in \GL_2(\Zp),\ \begin{pmatrix} a & b \\ c & d \end{pmatrix} \equiv \begin{pmatrix} 1 & 0 \\ 0 & 1 \end{pmatrix} \bmod p^{r} \right\}
\]
以及 $K_1(M;p^{r}) := K_1^{p}(M)K_p(r)$.
它们分别是 $\GL_2(\widehat{\mathbb{Z}}{}^{p})$, $\GL_2(\Zp)$ 和 $\GL_2(\widehat{\mathbb{Z}})$ 的开子群.
记 $Y_1(M;p^{r}) := Y(K_1(M;p^{r}))$ 为分类带 $p$ 之外 $\Gamma_1(M)$ 级结构、$p$ 处 $\Gamma(p^{r})$ 级结构的椭圆曲线的 (开) 模曲线.

$\Zp$-模 $X$ 的 $p$ 进完备化按定义是 $\Zp$-模:
\[
\widehat{X} := \varprojlim_n X/p^{n}X.
\]
对 $i \in \{0,1\}$ (resp.\ $i \in \{1,2\}$), 记 $\widehat{H}{}^{i}(K_1^{p}(M))$ (resp.\ $\widehat{H}{}_{c}^{i}(K_1^{p}(M))$) 为 $\Zp$-模 $\varinjlim_r H^{i}(Y_1(M;p^{r}),\Zp)$ (resp.\ $\Zp$-模 $\varinjlim_r H_{c}^{i}(Y_1(M;p^{r}),\Zp)$) 的 $p$ 进完备化, 并对 $\Qp$ 的任何有限扩张 $L$ 记 $\widehat{H}{}^{i}(K_1^{p}(M))_L := \widehat{H}{}^{i}(K_1^{p}(M)) \otimes_{\Zp} L$ (resp.\ $\widehat{H}{}_{c}^{i}(K_1^{p}(M))_L := \widehat{H}{}_{c}^{i}(K_1^{p}(M)) \otimes_{\Zp} L$).
$\Zp$-模 $\widehat{H}{}^{i} = \widehat{H}{}^{i}(K_1^{p}(M))$ 和 $\widehat{H}{}_{c}^{i} = \widehat{H}{}_{c}^{i}(K_1^{p}(M))$ 无挠, 并配备由其在 $\varinjlim_r H^{i}(Y_1(M;p^{r}),\Zp)$ 和 $\varinjlim_r H_{c}^{i}(Y_1(M;p^{r}),\Zp)$ 上的作用诱导的 $\Gal(\Q/\Q)$, $\GL_2(\Qp)$ 和 $\mathbb{T}$ 的自然作用.
$\GL_2(\Qp)$ 的作用使 $\widehat{H}{}^{i}(K_1^{p}(M))_L$ 和 $\widehat{H}{}_{c}^{i}(K_1^{p}(M))_L$ 成为可容许的酉 $\GL_2(\Qp)$-Banach 空间 (\cite[Th.~2.2.13]{ref:24}).
若 $i = 1$, 由 \cite[Lem.~2.2.1]{ref:9} 或 \cite[Prop.~4.3.9]{ref:24}, 从紧支上同调到上同调的自然映射:
\begin{equation}
\label{eq:7}
\widehat{H}{}_{c}^{1}(K_1^{p}(M)) \to \widehat{H}{}^{1}(K_1^{p}(M))
\end{equation}
是满射, 记 $\widehat{M}$ 为其核.
它仍配备 $\Gal(\Q/\Q)$, $\GL_2(\Qp)$ 和 $\mathbb{T}$ 的诱导作用.
这个核有一个简单的显式描述, 我们下面给出.

若 $K_f$ 是 $\GL_2(\mathbb{A}_f)$ 的任何紧开子群, 回顾模曲线 $Y(K_f)$ 的尖点集合有 adele 描述:
\[
C(K_f) := \GL_2(\Q)\backslash\bigl(\mathbb{P}^{1}(\Q) \times \pi_0 \times \GL_2(\mathbb{A}_f)\bigr)/K_f \xrightarrow{\sim} T(\Q)^{+}N(\mathbb{A}_f)\backslash \GL_2(\mathbb{A}_f)/K_f,
\]
其中 $T(\Q)^{+}$ 是 $T(\Q)$ 中正行列式对角矩阵的子群, $\pi_0$ 是 $\GL_2(\mathbb{R})$ 的连通分支群.
对 $K_f$ 取射影极限, 得到 $\widehat{C} := \varprojlim_{K_f} C(K_f)$ 这个 profinite 集合的如下描述:
\begin{equation}
\label{eq:8}
\widehat{C} \xrightarrow{\sim} T(\Q)^{+}N(\mathbb{A}_f)\backslash \GL_2(\mathbb{A}_f) \xrightarrow{\sim} (\pm N(\widehat{\mathbb{Z}}))\backslash \GL_2(\widehat{\mathbb{Z}}).
\end{equation}
"行列式"映射诱导一个满射:
\begin{equation}
\label{eq:9}
\widehat{C} \twoheadrightarrow \Q_{>0}^{\times}\backslash \mathbb{A}_f^{\times} \xrightarrow{\sim} \widehat{\mathbb{Z}}^{\times},
\end{equation}
它等同于把尖点送到包含它的连通分支的满射 $C(K_f) \to \pi_0(Y(K_f))$ 的射影极限.
由于 $T(\mathbb{A}_f)$ 正规化 $T(\Q)^{+}N(\mathbb{A}_f)$, 用 $T(\mathbb{A}_f)$ 左乘 (经 \eqref{eq:8}) 诱导 $T(\mathbb{A}_f)$ 在 $\widehat{C}$ 上的作用.
定义一个同态 $\Gal(\Q/\Q) \to T(\mathbb{A}_f)$:
\[
\sigma \mapsto \begin{pmatrix} \mathrm{cyclo}(\sigma) & 0 \\ 0 & 1 \end{pmatrix},
\]
其中 $\sigma \in \Gal(\Q/\Q)$, $\mathrm{cyclo} : \Gal(\Q/\Q) \to \widehat{\mathbb{Z}}^{\times}$ 是 adele 分圆特征标.
这个同态与 $T(\mathbb{A}_f)$ 在 $\widehat{C}$ 上的作用诱导 $\Gal(\Q/\Q)$ 在 $\widehat{C}$ 上的一个作用, 它等同于 $\Gal(\Q/\Q)$ 在 $C(K_f)$ 上自然作用的射影极限.

记 $\mathcal{C}^{0,\mathrm{lc}}(\widehat{C},\Zp)$ (resp.\ $\mathcal{C}^{0,\mathrm{lc}}(\widehat{\mathbb{Z}}^{\times},\Zp)$) 为从 $\widehat{C}$ (resp.\ $\widehat{\mathbb{Z}}^{\times}$) 到 $\Zp$ 的在 $p$ 进变量处连续、在其它变量处局部常值的函数空间.
$\GL_2(\mathbb{A}_f)$ 在 $\widehat{C}$ 上的右作用使 $\mathcal{C}^{0,\mathrm{lc}}(\widehat{C},\Zp)$ 成为 $\GL_2(\mathbb{A}_f)$-表示.
经 \eqref{eq:9}, $\mathcal{C}^{0,\mathrm{lc}}(\widehat{\mathbb{Z}}^{\times},\Zp)$ 等同于 $\mathcal{C}^{0,\mathrm{lc}}(\widehat{C},\Zp)$ 的一个不变子 $\Zp$-模, 并设:
\[
\widetilde{\mathcal{C}}{}^{0,\mathrm{lc}}(\widehat{C},\Zp) := \mathcal{C}^{0,\mathrm{lc}}(\widehat{C},\Zp)/\mathcal{C}^{0,\mathrm{lc}}(\widehat{\mathbb{Z}}^{\times},\Zp).
\]
它是无挠 $\Zp$-模, 配备 $T(\mathbb{A}_f)$ 和 $\Gal(\Q/\Q)$ 经其在 $\mathcal{C}^{0,\mathrm{lc}}(\widehat{C},\Zp)$ 上的作用 (后者由 $T(\mathbb{A}_f)$ 在 $\widehat{C}$ 上的作用推出) 诱导的作用.
这些作用与 $\GL_2(\mathbb{A}_f)$ 的作用交换.

\begin{proposition}\label{prop:3.1.1}
配备 $\Gal(\Q/\Q)$, $\GL_2(\Qp)$ 和 $\mathbb{T}$ 作用的 $\Zp$-模 $\widehat{M}$ 同构于 $\widetilde{\mathcal{C}}{}^{0,\mathrm{lc}}(\widehat{C},\Zp)$ 的子 $\Zp$-模 $\widetilde{\mathcal{C}}{}^{0,\mathrm{lc}}(\widehat{C},\Zp)^{K_1^{p}(M)}$ 及其上面诱导的 $\Gal(\Q/\Q)$ 和 $\GL_2(\Qp)$ 作用, 其中对 $\ell \nmid Mp$, $T_{\ell}$ (resp.\ $S_{\ell}$) 的作用是 $\GL_2(\Zp)$ 在 $\GL_2(\Qp)$ 中的 Hecke 代数中双陪集 $\GL_2(\mathbb{Z}_{\ell})\begin{pmatrix} \ell & 0 \\ 0 & 1 \end{pmatrix}\GL_2(\mathbb{Z}_{\ell})$ (resp.\ $\GL_2(\mathbb{Z}_{\ell})\begin{pmatrix} \ell & 0 \\ 0 & \ell \end{pmatrix}\GL_2(\mathbb{Z}_{\ell})$) 的作用.
\end{proposition}

\begin{proof}
这是众所周知的结果, 例如对 \cite[\S2.1]{ref:9} 的正合列 (1) 关于 $n$ 取射影极限得到.
关于 $\mathbb{T}$ 在 $\widehat{M}$ 上的经典作用与 (在素数 $\ell$, $(\ell,M) = 1$ 处) Hecke 代数在 $\widetilde{\mathcal{C}}{}^{0,\mathrm{lc}}(\widehat{C},\Zp)^{K_1^{p}(M)}$ 上的作用的比较, 参见 \cite[Prop.~4.4.2]{ref:24} 的证明.
\end{proof}

有必要显式描述 $\mathbb{T}$ 在 $\widetilde{\mathcal{C}}{}^{0,\mathrm{lc}}(\widehat{C},\Zp)^{K_1^{p}(M)}$ 上的作用.
若 $a \in \mathbb{Z}_{Mp}^{\times}$, 记 $a_p \in \Zp^{\times}$ 为 $a$ 经自然投影 $\mathbb{Z}_{Mp}^{\times} \to \Zp^{\times}$ 的像, $\widehat{a} \in \widehat{\mathbb{Z}}^{\times}$ 为经自然投影 $\widehat{\mathbb{Z}}^{\times} \to \mathbb{Z}_{Mp}^{\times}$ 送到 $a$ 的 $\widehat{\mathbb{Z}}^{\times}$ 的任一元素.
由于 $T(\mathbb{A}_f)$ 在 $\mathcal{C}^{0,\mathrm{lc}}(\widehat{C},\Zp)$ 上的作用与 $\GL_2(\mathbb{A}_f)$ 的作用交换, 它保持子模 $\mathcal{C}^{0,\mathrm{lc}}(\widehat{C},\Zp)^{K_1^{p}(M)}$.
我们将用 $T(\mathbb{A}_f)$ 的这个作用在 $\mathcal{C}^{0,\mathrm{lc}}(\widehat{C},\Zp)^{K_1^{p}(M)}$ 上定义 $\mathbb{Z}_{Mp}^{\times}$ 的两个作用.
对 $a \in \mathbb{Z}_{Mp}^{\times}$, $f \in \mathcal{C}^{0,\mathrm{lc}}(\widehat{C},\Zp)^{K_1^{p}(M)}$ 和 $\widehat{c} \in \widehat{C}$ 设:
\[
(\langle a \rangle_1 f)(\widehat{c}) := a_p f\!\left(\begin{pmatrix} \widehat{a}^{-1} & 0 \\ 0 & 1 \end{pmatrix}\widehat{c}\right)
\]
(这里 $a_p$ 简单经由乘以 $a_p \in \Zp^{\times}$ 作用) 以及:
\[
(\langle a \rangle_2 f)(\widehat{c}) := f\!\left(\begin{pmatrix} 1 & 0 \\ 0 & \widehat{a}^{-1} \end{pmatrix}\widehat{c}\right).
\]
利用如下等同可以验证这些作用良定义, 与提升 $\widehat{a}$ 的选择无关:
\begin{equation}
\label{eq:10}
\mathcal{C}^{0,\mathrm{lc}}(\widehat{C},\Zp)^{K_1^{p}(M)} = \mathcal{C}^{0}(\widehat{C}/K_1^{p}(M),\Zp),
\end{equation}
其中右边的模是 $\widehat{C}/K_1^{p}(M) \xrightarrow{\sim} T(\Q)^{+}N(\mathbb{A}_f)\backslash \GL_2(\mathbb{A}_f)/K_1^{p}(M)$ 上取值于 $\Zp$ 的连续函数模.
这两个作用都保持 $\mathcal{C}^{0,\mathrm{lc}}(\widehat{C},\Zp)^{K_1^{p}(M)}$ 的子模 $\mathcal{C}^{0,\mathrm{lc}}(\widehat{\mathbb{Z}}^{\times},\Zp)^{K_1^{p}(M)}$, 从而在 $\widetilde{\mathcal{C}}{}^{0,\mathrm{lc}}(\widehat{C},\Zp)^{K_1^{p}(M)}$ 上诱导 (记法相同的) $\mathbb{Z}_{Mp}^{\times}$ 作用.

\begin{lemma}\label{lem:3.1.2}
若 $\ell$ 不整除 $Mp$, 则 $T_{\ell}$ (resp.\ $\ell S_{\ell}$) 在 $\widetilde{\mathcal{C}}{}^{0,\mathrm{lc}}(\widehat{C},\Zp)^{K_1^{p}(M)}$ 上的作用由 $\langle \ell \rangle_1 + \langle \ell \rangle_2$ (resp.\ $\langle \ell \rangle_1\langle \ell \rangle_2$) 给出.
\end{lemma}

\begin{proof}
利用 \eqref{eq:10} 和 \eqref{eq:8} 给出的 $\widehat{C}$ 的 adele 描述做简单计算即可.
\end{proof}

\begin{corollary}\label{cor:3.1.3}
Hecke 模 $\widehat{M}$ 是 Eisenstein 的.
\end{corollary}

\begin{proof}
设 $L$ 是 $\Qp$ 中 $\Qp$ 的有限扩张, $\lambda$ 是使 $\widehat{M}_L^{\lambda} \neq 0$ 的 Hecke 特征值系.
由命题 \ref{prop:3.1.1}, 可以把 $\mathbb{Z}_{Mp}^{\times}$ 在 $\widetilde{\mathcal{C}}{}^{0,\mathrm{lc}}(\widehat{C},\Zp)^{K_1^{p}(M)}$ 上的两个作用 $\langle\ \rangle_1$ 和 $\langle\ \rangle_2$ 转移到 $\widehat{M}$ 上, 从而转移到 $\widehat{M} \otimes_{\Zp} L$ 上.
这些作用保持 $\widehat{M}_L^{\lambda} \neq 0$, 并由引理 \ref{lem:3.1.2} 满足等式:
\begin{equation}
\label{eq:11}
\langle \ell \rangle_1 + \langle \ell \rangle_2 = \lambda(T_{\ell}), \qquad \langle \ell \rangle_1\langle \ell \rangle_2 = \ell\lambda(S_{\ell}).
\end{equation}
$\widehat{M}_L^{\lambda}$ 的自同构 $\langle \ell \rangle_i$ ($i = 1,2$) 因此满足 $\langle \ell \rangle_i^{2} - \lambda(\ell)\langle \ell \rangle_i + \ell\lambda(S_{\ell}) = 0$.
必要时扩大 $L$, 可以在 $L$ 上分解这些二次方程 (对一切 $\ell$), 并得出结论: $\widehat{M}_L^{\lambda}$ 包含一个非零子空间, 在其上 $\langle\ \rangle_1$ 和 $\langle\ \rangle_2$ 经 $\mathbb{Z}_{Mp}^{\times}$ 的一个特征标作用.
此外, 这些特征标必然是连续的, 因为自同构 $\langle\ \rangle_1$ 和 $\langle\ \rangle_2$ 就是连续的.
由 \eqref{eq:11}, 这个特征值系确实是 Eisenstein 的.
\end{proof}

\subsection{与模形式 $f$ 相关联的 $p$ 进 $\GL_2(\Qp)$-Banach 空间 $\Pi_p(f)$ 的定义}\label{sec:3.2}

对一切 $k \ge 2$ 和一切如 \ref{sec:3.1} 的 $(M,r)$, 记 $V_{k-2}$ 为 $Y_1(M;p^{r})$ 上与 $\GL_2(\Q)$ 的表示 $\Sym^{k-2}\Qp^{2}$ 相关联的局部系, 并设:
\[
H^{1}(K_1^{p}(M),V_{k-2}) := \varinjlim_r H^{1}(Y_1(M;p^{r}),V_{k-2}), \qquad
H_{c}^{1}(K_1^{p}(M),V_{k-2}) := \varinjlim_r H_{c}^{1}(Y_1(M;p^{r}),V_{k-2}).
\]
这些空间自然配备互相交换的 $\Gal(\Q/\Q)$, $\GL_2(\Qp)$ 和 $\mathbb{T}$ 作用.
此外, $\GL_2(\Qp)$ 的作用是光滑的.
由 \cite[\S4.3]{ref:24}, 对 $\Qp$ 的任何有限扩张 $L$, 有与 $\GL_2(\Qp)$ 和 $\mathbb{T}$ 交换的注入:
\begin{equation}
\label{eq:12}
(\Sym^{k-2}L^{2})^{\vee} \otimes_{\Qp} H_{*}^{1}(K_1^{p}(M),V_{k-2}) \hookrightarrow \widehat{H}{}_{*}^{1}(K_1^{p}(M))_L,
\end{equation}
(其中 $* \in \{\emptyset,c\}$), 其像恰好是 $\widehat{H}{}_{*}^{1}(K_1^{p}(M))_L$ 的 $(\Sym^{k-2}L^{2})^{\vee}$-局部代数向量子空间.

设 $f$ 是权 $k \ge 2$、水平 $N$、特征标 $\chi : (\mathbb{Z}/N\mathbb{Z})^{\times} \to \overline{\Q}^{\times}$ 的归一化尖点模形式, 定义在 $\Qp$ 中 $\Qp$ 的有限扩张 $L$ 上.
假设 $f$ 是 Hecke 算子 $T_{\ell}$ ($\ell \nmid N$) 的特征向量, 记 $a_{\ell} \in L$ 为相应特征值.
若 $N = Mp^{r}$ ($r \ge 0$, $(M,p) = 1$), 由 Eichler--Shimura 有 $\Gal(\Q/\Q)$-与 $\GL_2(\Qp)$-等变同构:
\begin{equation}
\label{eq:13}
\sigma(f) \otimes_L \pi_p(f) \otimes_L \pi^{p}(f)^{K_1^{p}(M)} \xrightarrow{\sim} H_{*}^{1}(K_1^{p}(M),V_{k-2})^{f},
\end{equation}
其中 $* \in \{\emptyset,c\}$, 非零 $L$-向量空间 $\pi^{p}(f)^{K_1^{p}(M)}$ 恰在 $f$ 在 $M$ 处为新时维数为 $1$.
结合 \eqref{eq:13} 与 \eqref{eq:12}, 得到与 $\Gal(\Q/\Q)$ 和 $\GL_2(\Qp)$ 作用相容的注入:
\begin{equation}
\label{eq:14}
\sigma(f) \otimes_L (\Sym^{k-2}L^{2})^{\vee} \otimes_L \pi_p(f) \otimes_L \pi^{p}(f)^{K_1^{p}(M)} \hookrightarrow \widehat{H}{}_{*}^{1}(K_1^{p}(M))^{f}.
\end{equation}
假设 $f$ 在 $M$ 处为新, 记 $\widehat{\pi}_p(f)$ 为局部代数表示 $(\Sym^{k-2}L^{2})^{\vee} \otimes_L \pi_p(f)$ 经 \eqref{eq:14} 给出的一个嵌入在 $\widehat{H}{}_{c}^{1}(K_1^{p}(M))_L$ 中的闭包 (它与嵌入无关), 并设:
\[
\Pi_p(f) := \Hom_{\Gal(\Q/\Q)}\bigl(\sigma(f),\widehat{H}{}_{c}^{1}(K_1^{p}(M))_L\bigr),
\]
给它配备由 $\widehat{H}{}_{c}^{1}(K_1^{p}(M))_L$ 上的作用诱导的 $\GL_2(\Qp)$ 作用.
两个 $L$-向量空间 $\widehat{\pi}_p(f)$ 和 $\Pi_p(f)$ 都是可容许的酉 $\GL_2(\Qp)$-Banach 空间, 且有 $\GL_2(\Qp)$-等变闭浸入 $\widehat{\pi}_p(f) \hookrightarrow \Pi_p(f)$.
不难验证 $\widehat{\pi}_p(f)$ 也是 $(\Sym^{k-2}L^{2})^{\vee} \otimes_L \pi_p(f)$ 在 $\widehat{H}{}^{1}(K_1^{p}(M))_L$ 中的闭包 (\cite[Prop.~2.2.3]{ref:9}).
证明 $\widehat{H}{}^{1}(K_1^{p}(M))_L$ 的 $\sigma(f)$-等型分量仍给出 $\Pi_p(f)$ 需要一些预备知识.

\begin{lemma}\label{lem:3.2.1}
不整除 $Mp$ 且使得 $\Frob_{\ell}$ 在 $\sigma(f)$ 下的像是纯量的素数 $\ell$ 的集合密度为零.
\end{lemma}

\begin{proof}
若 $P$ 表示所有素数的集合, $S \subseteq P$ 为任意子集, 回顾 $S$ 的密度 (若存在) 按定义是:
\[
d(S) := \lim_{x \to +\infty} \frac{\#\{\ell \in S \mid \ell \le x\}}{\#\{\ell \in P \mid \ell \le x\}} \le 1.
\]
设 $\sigma(f)^{0} : \Gal(\Q/\Q) \to \GL_2(\mathcal{O}_L)$ 是 $\sigma(f)$ 中 $\Gal(\Q/\Q)$-稳定的 $\mathcal{O}_L$-格, $S$ (resp.\ 对 $n \in \mathbb{Z}_{>0}$ 的 $S_n$) 为不整除 $Mp$ 且 $\Frob_{\ell}$ 在 $\sigma(f)$ (resp.\ $\sigma(f)^{0} \otimes_{\mathcal{O}_L} \mathcal{O}_L/p^{n}\mathcal{O}_L$) 下的像是纯量的素数 $\ell$ 的集合.
记 $G_n \subseteq \mathrm{PGL}_2(\mathcal{O}_L/p^{n}\mathcal{O}_L)$ 为 $\Gal(\Q/\Q)$ 经射影表示:
\[
\Gal(\Q/\Q) \xrightarrow{\sigma(f)^{0}} \GL_2(\mathcal{O}_L) \twoheadrightarrow \GL_2(\mathcal{O}_L/p^{n}\mathcal{O}_L) \twoheadrightarrow \mathrm{PGL}_2(\mathcal{O}_L/p^{n}\mathcal{O}_L)
\]
的像, $K_n \subset \overline{\Q}$ 为满足 $\Gal(K_n/\Q) \xrightarrow{\sim} G_n$ 的 $\Q$ 的有限 Galois 扩张.
对扩张 $K_n/\Q$ 应用 \v{C}ebotarev 定理 (参见 e.g.\ \cite[\S2.1]{ref:46}), 得到 $d(S_n)$ 存在且等于:
\[
d(S_n) = \frac{1}{\#\Gal(K_n/\Q)} = \frac{1}{\#G_n}.
\]
但由 \cite[Th.~4.3]{ref:39}, $\Gal(\Q/\Q)$ 经表示 $\sigma(f)^{0}$ "射影化"在 $\mathrm{PGL}_2(\mathcal{O}_L)$ 中的像是无限的, 故 $G_n$ 的基数随 $n$ 趋于 $+\infty$, 从而:
\begin{equation}
\label{eq:15}
\lim_{n \to +\infty} d(S_n) = 0.
\end{equation}
但对一切 $n$ 有:
\[
0 \le \limsup_{x \to +\infty} \frac{\#\{\ell \in S \mid \ell \le x\}}{\#\{\ell \in P \mid \ell \le x\}} \le \limsup_{x \to +\infty} \frac{\#\{\ell \in S_n \mid \ell \le x\}}{\#\{\ell \in P \mid \ell \le x\}} = d(S_n),
\]
故用 \eqref{eq:15} 令 $n \to +\infty$:
\[
\limsup_{x \to +\infty} \frac{\#\{\ell \in S \mid \ell \le x\}}{\#\{\ell \in P \mid \ell \le x\}} = 0 = \lim_{x \to +\infty} \frac{\#\{\ell \in S \mid \ell \le x\}}{\#\{\ell \in P \mid \ell \le x\}} = d(S),
\]
这就完成了证明.
\end{proof}

若 $S$ 是不整除 $Mp$ 的素数的集合, $X$ 是 Hecke 模, 记:
\[
X_{L}^{S,f} := \{x \in X \otimes_{\Zp} L \mid \forall\, \ell \in S,\ T_{\ell}x = a_{\ell}x,\ S_{\ell}x = \ell\chi(\ell)x\}.
\]

\begin{lemma}\label{lem:3.2.2}
设 $S$ 是不整除 $Mp$ 且密度为 $1$ 的素数集合, 则注入 $\widehat{H}{}_{c}^{1}(K_1^{p}(M))_L^{S,f} \hookrightarrow \widehat{H}{}_{c}^{1}(K_1^{p}(M))_L^{f}$ 和 $\widehat{H}{}^{1}(K_1^{p}(M))_L^{S,f} \hookrightarrow \widehat{H}{}^{1}(K_1^{p}(M))_L^{f}$ 是同构.
\end{lemma}

\begin{proof}
记 $(\widehat{H}{}_{c}^{1}(K_1^{p}(M)) \otimes_{\Zp} \mathcal{O}_L)^{S,f} := (\widehat{H}{}_{c}^{1}(K_1^{p}(M)) \otimes_{\Zp} \mathcal{O}_L) \cap \widehat{H}{}_{c}^{1}(K_1^{p}(M))_L^{S,f}$, 并设 $x = (x_n) \in (\widehat{H}{}_{c}^{1}(K_1^{p}(M)) \otimes_{\Zp} \mathcal{O}_L)$, 其中:
\[
x_n \in H_{c}^{1}(K_1^{p}(M),\mathcal{O}_L/p^{n}\mathcal{O}_L) := \varinjlim_m H_{c}^{1}(Y_1(M;p^{m}),\mathcal{O}_L/p^{n}\mathcal{O}_L).
\]
只需证明: 对一切 $n$, 对一切不整除 $Mp$ 的素数 $\ell$, $T_{\ell}x_n = a_{\ell}x_n$ 且 $S_{\ell}x_n = \ell\chi(\ell)x_n$.
取 $m$ 使得 $x_n \in H_{c}^{1}(Y_1(M;p^{m}),\mathcal{O}_L/p^{n}\mathcal{O}_L)$, $\ell$ 为不整除 $Mp$ 的素数, $\widehat{x}_n$ 为 $x_n$ 在 $H_{c}^{1}(Y_1(M;p^{m}),\mathcal{O}_L)$ 中的一个提升, 并取 $S$ 中的素数序列 $(\ell_i)_{i \in \mathbb{N}}$ 使得在 $\Q$ 的在 $Mp$ 之外非分歧的最大扩张的 Galois 群中有 $\Frob_{\ell_i} \to \Frob_{\ell}$ (当 $i \to +\infty$) (注意这意味着 $\ell_i$ 在 $\Zp^{\times}$ 中趋于 $\ell$).
由于 $S$ 密度为 $1$, 这样的序列由 \v{C}ebotarev 定理保证存在.
由 Eichler--Shimura 关系 (以及 $H_{c}^{1}(Y_1(M;p^{m}),\mathcal{O}_L) \otimes L$ 的已知 Galois 结构), 作为 $H_{c}^{1}(Y_1(M;p^{m}),\mathcal{O}_L)$ 的自同态还有 $T_{\ell_i} \to T_{\ell}$ 和 $S_{\ell_i} \to S_{\ell}$.
特别地, 对 $i \gg 0$, 有 $T_{\ell_i}x_n = a_{\ell_i}x_n = a_{\ell}x_n$ 和 $T_{\ell_i}x_n = T_{\ell}x_n$, 故 $T_{\ell}x_n = a_{\ell}x_n$.
同样, 对 $i \gg 0$, 有 $S_{\ell_i}x_n = \ell_i\chi(\ell_i)x_n = \ell\chi(\ell)x_n$ 和 $S_{\ell_i}x_n = S_{\ell}x_n$, 故 $S_{\ell}x_n = \ell\chi(\ell)x_n$.
对 $\widehat{H}{}^{1}(K_1^{p}(M))_L$ 的证明相同.
\end{proof}

\begin{proposition}\label{prop:3.2.3}
有同构:
\[
\Hom_{\Gal(\Q/\Q)}\bigl(\sigma(f),\widehat{H}{}_{c}^{1}(K_1^{p}(M))_L^{f}\bigr) \xrightarrow{\sim} \Pi_p(f).
\]
\end{proposition}

\begin{proof}
记 $S$ 为不整除 $Mp$ 且 $\Frob_{\ell}$ 在 $\sigma(f)$ 上不是纯量的素数集合.
由引理 \ref{lem:3.2.1}, $S$ 密度为 $1$, 且由引理 \ref{lem:3.2.2}, 只需用 $\widehat{H}{}_{c}^{1}(K_1^{p}(M))_L^{S,f}$ 代替 $\widehat{H}{}_{c}^{1}(K_1^{p}(M))_L^{f}$ 证明该命题.
给 $\Pi_p(f)$ 配备由 Hecke 算子在 $\widehat{H}{}_{c}^{1}(K_1^{p}(M))_L$ 上的作用 (它与 $\Gal(\Q/\Q)$ 的作用交换) 诱导的作用.
只需证明: 对一切 $x \in \Pi_p(f)$ 和一切 $\ell \in S$, 有 $T_{\ell}(x) = a_{\ell}x$ 和 $S_{\ell}(x) = \ell^{k-2}\chi(\ell)x$.
有规范 Galois-与 Hecke-等变注入:
\[
\sigma(f) \otimes_{\mathcal{O}_L} \Pi_p(f) \hookrightarrow \widehat{H}{}_{c}^{1}(K_1^{p}(M))_L,\qquad v \otimes x \mapsto x(v)
\]
(注入由 $\sigma(f)$ 的不可约性得到).
取 $(v,x) \in \sigma(f) \times \Pi_p(f)$, 则 Eichler--Shimura 关系 $\Frob_{\ell}^{-2} - T_{\ell} \circ \Frob_{\ell}^{-1} + \ell S_{\ell} = 0$ (通过取有限水平的这些关系的极限, 它们在 $\widehat{H}{}_{c}^{1}(K_1^{p}(M))$ 上仍成立) 与 $\sigma(f)|_{\Gal(\Q_{\ell}/\Q_{\ell})}$ 的性质蕴含在 $\sigma(f) \otimes_L \Pi_p(f)$ 中, 对一切 $v \in \sigma(f)$ 和一切 $x \in \Pi_p(f)$, 有等式:
\[
-\Frob_{\ell}^{-1}(v) \otimes (T_{\ell} - a_{\ell})(x) + v \otimes (\ell S_{\ell} - \ell^{k-1}\chi(\ell))(x) = 0.
\]
由于 $\Frob_{\ell}$ 在 $\sigma(f)$ 上不是纯量, 固定 $x$ 并让 $v$ 在 $\sigma(f)$ 中变化, 容易推出 $(T_{\ell} - a_{\ell})(x) = 0$ 和 $(S_{\ell} - \ell^{k-2}\chi(\ell))(x) = 0$.
\end{proof}

\begin{proposition}\label{prop:3.2.4}
有同构:
\[
\Pi_p(f) \xrightarrow{\sim} \Hom_{\Gal(\Q/\Q)}\bigl(\sigma(f),\widehat{H}{}^{1}(K_1^{p}(M))_L^{f}\bigr)
\xrightarrow{\sim} \Hom_{\Gal(\Q/\Q)}\bigl(\sigma(f),\widehat{H}{}^{1}(K_1^{p}(M))_L\bigr).
\]
\end{proposition}

\begin{proof}
左边的同构借助引理 \ref{lem:3.2.1} 和 \ref{lem:3.2.2} 完全像命题 \ref{prop:3.2.3} 那样证明.
因此只需证明自然映射:
\begin{equation}
\label{eq:16}
\Hom_{\Gal(\Q/\Q)}\bigl(\sigma(f),(\widehat{H}{}_{c}^{1}(K_1^{p}(M))_L)^{f}\bigr) \to \Hom_{\Gal(\Q/\Q)}\bigl(\sigma(f),(\widehat{H}{}^{1}(K_1^{p}(M))_L)^{f}\bigr)
\end{equation}
是同构.
记 $\mathbb{T}_{\mathcal{O}_L} := \mathbb{T} \otimes_{\Zp} \mathcal{O}_L$, $I \subset \mathbb{T}_{\mathcal{O}_L}$ 为由 $T_{\ell} - a_{\ell}$ 和 $S_{\ell} - \ell^{k-2}\chi(\ell)$ 生成的理想.
满射 \eqref{eq:7} 诱导正合列:
\[
0 \to \widehat{M}_L \to \widehat{H}{}_{c}^{1}(K_1^{p}(M))_L \to \widehat{H}{}^{1}(K_1^{p}(M))_L \to 0.
\]
应用函子 $\Hom_{\mathbb{T}_{\mathcal{O}_L}}(\mathbb{T}_{\mathcal{O}_L}/I,\ \cdot)$, 得到正合列:
\[
0 \to (\widehat{M}_L)^{f} \to (\widehat{H}{}_{c}^{1}(K_1^{p}(M))_L)^{f} \to (\widehat{H}{}^{1}(K_1^{p}(M))_L)^{f} \to \Ext_{\mathbb{T}_{\mathcal{O}_L}}^{1}(\mathbb{T}_{\mathcal{O}_L}/I,\widehat{M}_L).
\]
由推论 \ref{cor:3.1.3}, 因 $f$ 是尖点的, 有 $(\widehat{M}_L)^{f} = 0$, 故有正合列:
\[
0 \to (\widehat{H}{}_{c}^{1}(K_1^{p}(M))_L)^{f} \to (\widehat{H}{}^{1}(K_1^{p}(M))_L)^{f} \to \Ext_{\mathbb{T}_{\mathcal{O}_L}}^{1}(\mathbb{T}_{\mathcal{O}_L}/I,\widehat{M}_L),
\]
它诱导新的正合列:
\[
0 \to \Hom_{\Gal(\Q/\Q)}\bigl(\sigma(f),(\widehat{H}{}_{c}^{1}(K_1^{p}(M))_L)^{f}\bigr)
\to \Hom_{\Gal(\Q/\Q)}\bigl(\sigma(f),(\widehat{H}{}^{1}(K_1^{p}(M))_L)^{f}\bigr)
\to \Hom_{\Gal(\Q/\Q)}\bigl(\sigma(f),\Ext_{\mathbb{T}_{\mathcal{O}_L}}^{1}(\mathbb{T}_{\mathcal{O}_L}/I,\widehat{M}_L)\bigr).
\]
由命题 \ref{prop:3.1.1}, $\Gal(\Q/\Q)$ 在 $\widehat{M}_L$ 上的作用经 $\Gal(\Q/\Q)^{\mathrm{ab}}$ 分解, 因此 $\Gal(\Q/\Q)$ 在 $\Ext_{\mathbb{T}_{\mathcal{O}_L}}^{1}(\mathbb{T}_{\mathcal{O}_L}/I,\widehat{M}_L)$ 上的作用也如此.
由于 $\sigma(f)$ 是绝对不可约的 $\Gal(\Q/\Q)$-表示, 由此推出:
\[
\Hom_{\Gal(\Q/\Q)}\bigl(\sigma(f),\Ext_{\mathbb{T}_{\mathcal{O}_L}}^{1}(\mathbb{T}_{\mathcal{O}_L}/I,\widehat{M}_L)\bigr) = 0,
\]
从而得到所求同构.
\end{proof}
